% =====================================================================
%  Test Book — English sample (verifies all template features)
%  Compile with XeLaTeX:
%      xelatex main.tex   (run 2x for the TOC)
%
%  Class options:
%    portuguese  — Portuguese labels (default: English)
%    notime      — hide compilation time on cover page
%    nochangelog — disable changelog rendering
%    noauthors   — hide authors on cover page
%    noanswers   — hide Remarks and Test Result rows in test cases
% =====================================================================
\documentclass{testbook}

% --- Document configuration -----------------------------------------
\settestbooktitle{POC Test Cases — Data Platform Validation}
\setheadertitle{Huawei Cloud -- POC Test Cases}
\setcovertext{Huawei Technologies CO., LTD}
\setdocversion{1.0.0}
\setdocdate{\today}
\setdocauthors{Wallace Wu}

\begin{document}

% ---- COVER ---------------------------------------------------------
\makecover

% ---- TABLE OF CONTENTS ---------------------------------------------
\maketoc

% ---- BODY (restarts page numbering at 1) ---------------------------
\startbody

\begin{infobox}
This document demonstrates the \inlinecode{testbook} template (\inlinecode{testbook.cls}) for
Huawei Cloud POC/acceptance test case documents. It exercises all available
commands and environments: cover page, table of contents, test case environment
with all field commands, objectives block, code blocks, images, callouts,
tables, badge, and changelog.
\end{infobox}

% =====================================================================
%  Section 1 — Test Cases Description
% =====================================================================
\section{Test Cases Description}

\begin{objectives}
  \generalobjective{Validate the Huawei Cloud Data Platform deployment for the
    POC environment, ensuring that all core services (DataArts Studio, MRS,
    OBS, VPC) are provisioned correctly, accessible, and functioning as
    specified in the solution architecture.}

  \prerequisites
  \begin{itemize}
    \item Huawei Cloud account with IAM admin privileges in the \inlinecode{sa-brazil-1} region.
    \item VPC, subnet, and security groups already created by the infrastructure team.
    \item DataArts Studio instance (Dayu) provisioned and licensed.
    \item MRS cluster deployed with the required component stack (HDFS, YARN, Spark, Hive).
    \item OBS bucket created and accessible with the designated IAM policy.
    \item Test execution environment with network connectivity to the Huawei Cloud Console.
  \end{itemize}
\end{objectives}

\subsection{Test Scope}

\begin{table}[H]
  \begin{hutable}{|l|p{0.72\textwidth}|}
    \rowcolor{huaweired} \thd{Domain} & \thd{What Should Be Observed} \\
    \tbody
    Platform Architecture & Core services are deployed, healthy, and reachable via the Console API. \\
    Data Engineering & DataArts Studio pipelines can ingest, transform, and load data into the target storage. \\
    Data Storage & OBS buckets and MRS HDFS accept writes/reads with correct permissions. \\
    Security & IAM policies, VPC isolation, and security group rules enforce the intended access control. \\
  \end{hutable}
  \caption{Test scope by domain.}
\end{table}

\subsection{Preconditions and Preparations}

\begin{itemize}
  \item All infrastructure resources must be provisioned via Terraform and the \inlinecode{terraform apply} output must show zero errors.
  \item The MRS cluster must be in \textbf{Running} state before any data engineering test case is executed.
  \item Test data files (CSV/JSON) must be uploaded to the designated OBS bucket prior to pipeline execution.
  \item Each tester must have a personal IAM account with the \textbf{Dayu\_User} role assigned in DataArts Studio.
  \item Network firewall rules must allow inbound TCP 22 (SSH) and TCP 443 (HTTPS) from the test execution subnet.
\end{itemize}

\subsection{Acceptance Method}

\begin{table}[H]
  \begin{hutable}{|l|p{0.72\textwidth}|}
    \rowcolor{huaweired} \thd{Result} & \thd{Explanation} \\
    \tbody
    Satisfied & The test case passes fully — observed behavior matches the expected result with no deviations. \\
    Satisfied but with reservations & The test case passes but minor deviations were observed that do not impact the POC objectives (documented in Remarks). \\
    Not satisfied & The test case fails — observed behavior does not match the expected result. A defect report must be filed. \\
    Untested & The test case could not be executed due to blocking dependencies or environment issues. \\
  \end{hutable}
  \caption{Acceptance criteria.}
\end{table}

% =====================================================================
%  Section 2 — Platform Architecture
% =====================================================================
\section{Platform Architecture}

\begin{testcase}{TC-001: Verify MRS cluster deployment and health}
  \testobjective{Confirm that the MapReduce Service (MRS) cluster is deployed
    with the required components (HDFS, YARN, Spark, Hive) and all nodes
    report a healthy status.}
  \testprerequisites{1. Terraform infrastructure apply completed successfully. \\ 2. MRS cluster ID is available in the Terraform output. \\ 3. IAM user has \textbf{MRS\_Viewer} or higher role.}
  \testprocedure{1. Log in to the Huawei Cloud Console. \\ 2. Navigate to \menu{Console, MapReduce Service, Clusters}. \\ 3. Locate the cluster by its name (e.g. \inlinecode{mrs-poc-cluster}). \\ 4. Verify the cluster status is \textbf{Running}. \\ 5. Click the cluster name and review the \textbf{Component} tab — confirm HDFS, YARN, Spark, and Hive are listed with status \textbf{Normal}. \\ 6. Review the \textbf{Node} tab — confirm all master and core nodes show status \textbf{Running}.}
  \testexpected{1. Cluster status is \textbf{Running}. \\ 2. All four components (HDFS, YARN, Spark, Hive) are listed with status \textbf{Normal}. \\ 3. All nodes report status \textbf{Running} with no alarms.}
  \testremarks{If any component shows \textbf{Abnormal}, check the MRS alarm page before marking the result.}
  \testresult{Pass}
\end{testcase}

\begin{testcase}{TC-002: Verify OBS bucket creation and accessibility}
  \testobjective{Confirm that the Object Storage Service (OBS) bucket used
    for the POC data lake is created, has the correct storage class, and
    is accessible with the designated IAM policy.}
  \testprerequisites{1. OBS bucket name is defined in the Terraform configuration. \\ 2. IAM policy granting read/write access to the bucket is attached to the test user.}
  \testprocedure{1. Log in to the Console. \\ 2. Navigate to \menu{Console, Object Storage Service}. \\ 3. Locate the bucket (e.g. \inlinecode{poc-datalake-raw}). \\ 4. Verify the storage class is \textbf{Standard}. \\ 5. Upload a test file (\inlinecode{test-upload.txt}) to the bucket. \\ 6. Download the file and compare its content with the original. \\ 7. Delete the test file.}
  \testexpected{1. Bucket exists with storage class \textbf{Standard}. \\ 2. File upload completes without error. \\ 3. Downloaded file content matches the original. \\ 4. File deletion succeeds.}
  \testremarks{OBS eventual consistency may cause a brief delay before a newly uploaded object appears in listings.}
  \testresult{Pass}
\end{testcase}

\begin{testcase}{TC-003: Verify VPC and security group configuration}
  \testobjective{Confirm that the VPC, subnet, and security groups are
    configured according to the network architecture diagram, with the
    correct CIDR blocks and inbound/outbound rules.}
  \testprerequisites{1. VPC and subnet IDs are available from the Terraform output. \\ 2. Security group names are documented in the solution architecture.}
  \testprocedure{1. Log in to the Console. \\ 2. Navigate to \menu{Console, Virtual Private Cloud, VPCs}. \\ 3. Locate the VPC (e.g. \inlinecode{vpc-poc}) and verify its CIDR block matches the architecture (e.g. \inlinecode{10.0.0.0/16}). \\ 4. Click the subnet and verify its CIDR (e.g. \inlinecode{10.0.1.0/24}). \\ 5. Navigate to \textbf{Security Groups} and verify inbound rules: TCP 22 from the bastion subnet, TCP 443 from the corporate proxy. \\ 6. Verify outbound rules allow all traffic (default).}
  \testexpected{1. VPC CIDR is \inlinecode{10.0.0.0/16}. \\ 2. Subnet CIDR is \inlinecode{10.0.1.0/24}. \\ 3. Inbound rules match the architecture: TCP 22 and TCP 443 from the specified sources. \\ 4. Outbound allows all traffic.}
  \testremarks{If the security group has additional rules beyond the architecture spec, document them in the Remarks field.}
  \testresult{Pass}
\end{testcase}

% =====================================================================
%  Section 3 — Data Engineering
% =====================================================================
\section{Data Engineering}

\begin{testcase}{TC-004: Verify DataArts Studio workspace creation}
  \testobjective{Confirm that the DataArts Studio (Dayu) instance is
    provisioned, the workspace is accessible, and the designated IAM
    users can log in with the \textbf{Dayu\_User} role.}
  \testprerequisites{1. DataArts Studio instance is provisioned and its status is \textbf{Running}. \\ 2. Test IAM user has the \textbf{Dayu\_User} role assigned.}
  \testprocedure{1. Log in to the Console with the test IAM user. \\ 2. Navigate to \menu{Console, DataArts Studio, Workspaces}. \\ 3. Locate the workspace (e.g. \inlinecode{poc-workspace}). \\ 4. Click \textbf{Access Workspace} to open the DataArts Studio console. \\ 5. Verify the left navigation panel loads with the expected modules: \textbf{Data Integration}, \textbf{Data Development}, \textbf{Data Architecture}.}
  \testexpected{1. Workspace is listed and accessible. \\ 2. DataArts Studio console opens without errors. \\ 3. All three modules (Data Integration, Data Development, Data Architecture) are visible in the navigation.}
  \testremarks{If the workspace fails to load, check the Dayu instance status and the IAM role assignment.}
  \testresult{Pass}
\end{testcase}

\begin{testcase}{TC-005: Verify data ingestion pipeline execution}
  \testobjective{Confirm that a DataArts Studio batch pipeline can
    successfully ingest CSV data from OBS, apply a basic transformation,
    and write the result to the target OBS path.}
  \testprerequisites{1. TC-004 completed (DataArts Studio workspace accessible). \\ 2. Source CSV file (\inlinecode{customers.csv}) exists in the OBS raw bucket. \\ 3. Target OBS path (\inlinecode{poc-datalake-curated/customers/}) is configured. \\ 4. Pipeline \inlinecode{pl-ingest-customers} is published in DataArts Studio.}
  \testprocedure{1. In the DataArts Studio console, navigate to \textbf{Data Development}. \\ 2. Locate the pipeline \inlinecode{pl-ingest-customers} in the job list. \\ 3. Click \textbf{Run} to execute the pipeline. \\ 4. Wait for the pipeline status to change to \textbf{Success} (timeout: 10 minutes). \\ 5. Navigate to OBS and verify the output file exists in \inlinecode{poc-datalake-curated/customers/}. \\ 6. Download the output file and verify the row count matches the source (no data loss).}
  \testexpected{1. Pipeline execution completes with status \textbf{Success}. \\ 2. Output file is created in the target OBS path. \\ 3. Row count of the output file matches the source file.}
  \testremarks{Pipeline execution time varies with data volume. For the POC, the source file contains approximately 10,000 rows.}
  \testresult{Pass}
\end{testcase}

\begin{testcase}{TC-006: Verify Spark SQL query on MRS cluster}
  \testobjective{Confirm that a Spark SQL query can be executed on the
    MRS cluster to read data from a Hive table and return correct
    results, validating the integration between MRS and the data lake.}
  \testprerequisites{1. TC-001 completed (MRS cluster healthy). \\ 2. Hive table \inlinecode{poc\_db.customers} exists and contains data loaded by the ingestion pipeline. \\ 3. SSH access to the MRS master node is configured.}
  \testprocedure{1. SSH into the MRS master node. \\ 2. Run the Spark SQL CLI: \inlinecode{spark-sql --master yarn -{}-conf spark.sql.hive.convertMetastoreOrc=true -e "SELECT COUNT(*) FROM poc\_db.customers;"} \\ 3. Verify the returned count matches the expected row count (10,000). \\ 4. Run a sample query to verify data integrity: \inlinecode{spark-sql --master yarn -e "SELECT customer\_id, name FROM poc\_db.customers LIMIT 5;"} \\ 5. Confirm the query returns 5 rows with non-null values.}
  \testexpected{1. \inlinecode{COUNT(*)} returns 10,000. \\ 2. The sample query returns 5 rows with valid \inlinecode{customer\_id} and \inlinecode{name} values.}
  \testremarks{If the Spark SQL CLI is not available, use the MRS Spark2x component's web UI to submit the query instead.}
  \testresult{Pass}
\end{testcase}

% =====================================================================
%  Changelog
% =====================================================================
\begin{changelog}
  \changelogentry{1.0.0}{\today}{
    \item Initial version.
    \item Added Platform Architecture test cases (TC-001 to TC-003).
    \item Added Data Engineering test cases (TC-004 to TC-006).
  }
\end{changelog}

\end{document}
